\documentclass[11pt, draft]{article}
\topmargin=0mm
\oddsidemargin=0mm
\textwidth170mm
\textheight230mm
\begin{document}
%\pagestyle{empty}
\language=0
%\hyphenation{}

\begin{center}
{\bf
APPLICATION OF COMPUTER ALGEBRA TO ANALYSIS OF DIFFERENTIAL
EQUATIONS SYSTEMS COMPATIBILITY
}\\
\vspace {0.3cm}
{S.V.~Meleshko, V.P.~Shapeev}\\
{\it \small ITAM SB RAS, 630090, Novosibirsk, Russia,
 e-mail: shapeew@itam.nsc.ru}
\end{center}
\vspace {0.6cm}

One of the fields of computer algebra application is
the theory of differential equations system compatibility.
Necessity of analysis on compatibility arises, for instance,
when performing group analysis of partial differential equations
(PDE), or applying the differential constraints method (DCM) [1].
According to Cartan theorem of compatibility theory,
after finite number of continuations
any system of differential equations can be
 reduced to involution (compatible system) or to contradiction.

At present, two equivalent compatibility analysis methods
 applicable to PDE exist: Cartan method and Janet-Spencer-Kuranishi
 (JSK) method.
 In applications, analysis on PDE compatibility often leads to
 huge symbolic computations, to which modern computer algebra systems (CAS)
 can be applied.
>From algorithmic point of view methods of analysis on nonlinear PDE systems
compatibility are nonlinear,
ramified algorithms, which are too complicated for CAS.
There are several attempts of theirs realizations.
 A special case
of reduction of one particular system was considered in [2], and a
principal ability of using a computer for solving such a problems
was shown.
We have used different CASs for realization of Cartan and JSK methods.
Full charts of algorithms in algebraic terms were written,
software inplementation in REFAL and REDUCE systems was made,
particular problems of reduction to involution some overdetermined PDE
systems, which
arise in hydroaerodynamics, were solved [1-7].
This work was partially supported by the Russian Foundation for Basic
Research (grants  99-01-00515 and 00-01-00370) and
INTAS (proposal 1222).


\begin{thebibliography}{}
\bibitem{si}
Sidorov A.F., Shapeev V.P., Yanenko N.N., {\it Technique
of differential conections and its applications in gas dynamics,}
Nauka, Novosibirsk, 1984, (Russian).
\bibitem{ch}
Shurygin V.A., Yanenko N.N., On computer realization of
algebraic/differential algorithms, in: {\it Problemy
Kibernetiki,} 6, 1961, (Russian).
\bibitem{ar}
  Arais E.A., Shapeev V.P. and Yanenko N.N. 1974 Computer realization
of Cartan's external forms, Dokl. AN USSR. {\bf~214}~(4),
737--738.
\bibitem{ga}
Ganzha V.G.,
 Meleshko C.V., Murzin F.A., Shapeev V.P., Yanenko N.N.
 Computer realization of the
algorithm for the investigation of compatibility of systems of
partial differential equations, {\it Doklady Akademii Nauk
USSR,} 261(5), 1044-1046, 1981, (Russian).
\bibitem{ga1}
V.G. Ganzha, S. V. Meleshko, V. P. Shapeev. Automation of
Compatibility Analysis of Quasilinear PDE Systems with the Aid
of the REDUSE System. Computer Algebra and Its Applications
to Mechanics. Nova Science Publishers, New York, 1993, p. 89-96.
\bibitem{me}
Meleshko S.V., On potential flows of polytropic gas of
the type of triple waves, {\it Prikladnaya Matematika i
Mekhanika,} 54(5), 775-779, 1990, (Russian).
\bibitem{me1}
S. V. Meleshko, V. P. Shapeev.
An application of the differential constraints method for
the two-dimensional equations of gas dynamics
{\it Prikladnaya Matematika i Mekhanika,}
 63(6), 909-916, 1999, (Russian).
\end{thebibliography}

\end{document}
